\documentclass[10pt,a4paper]{amsart}
\usepackage{amsmath,amssymb,amsthm,graphicx}
% Retained preamble for the Stacks Project chapter {CH} (shared project preamble).
% The chapter file pinned-original-source.tex only carries \input{{preamble}} on its first line,
% so no preamble text lives inside the chapter itself.  The Stacks Project uses one single shared
% preamble for every chapter: the snapshot stages it three times (sources/stacks-chow/preamble.tex,
% sources/stacks-etale-cohomology/preamble.tex and focused-seeds/AGP/sources/stacks-snapshot-20261011/
% preamble.tex) and those copies are identical in size.  The block below is copied mechanically,
% line for line, from sources/stacks-chow/preamble.tex lines 14-18 and 154-217: the Xy-pic setup and
% the native theorem declarations together with the mathematical macros.  Nothing was added, removed
% or redefined; the wrapper class and the amsmath/amssymb/amsthm/graphicx base packages are supplied
% by the assembler.
% For commutative diagrams we use Xy-pic
% Chapter: stacks-spaces-more-morphisms.tex -- More on Morphisms of Spaces.
\usepackage[all]{xy}

% We use 2cell for 2-commutative diagrams.
\xyoption{2cell}

% Theorem environments.
%
\theoremstyle{plain}
\newtheorem{theorem}[subsection]{Theorem}
\newtheorem{proposition}[subsection]{Proposition}
\newtheorem{lemma}[subsection]{Lemma}

\theoremstyle{definition}
\newtheorem{definition}[subsection]{Definition}
\newtheorem{example}[subsection]{Example}
\newtheorem{exercise}[subsection]{Exercise}
\newtheorem{situation}[subsection]{Situation}

\theoremstyle{remark}
\newtheorem{remark}[subsection]{Remark}
\newtheorem{remarks}[subsection]{Remarks}

\numberwithin{equation}{subsection}

% Macros
%
\def\lim{\mathop{\mathrm{lim}}\nolimits}
\def\colim{\mathop{\mathrm{colim}}\nolimits}
\def\Spec{\mathop{\mathrm{Spec}}}
\def\Hom{\mathop{\mathrm{Hom}}\nolimits}
\def\Ext{\mathop{\mathrm{Ext}}\nolimits}
\def\SheafHom{\mathop{\mathcal{H}\!\mathit{om}}\nolimits}
\def\SheafExt{\mathop{\mathcal{E}\!\mathit{xt}}\nolimits}
\def\Sch{\mathit{Sch}}
\def\Mor{\mathop{\mathrm{Mor}}\nolimits}
\def\Ob{\mathop{\mathrm{Ob}}\nolimits}
\def\Sh{\mathop{\mathit{Sh}}\nolimits}
\def\NL{\mathop{N\!L}\nolimits}
\def\CH{\mathop{\mathrm{CH}}\nolimits}
\def\proetale{{pro\text{-}\acute{e}tale}}
\def\etale{{\acute{e}tale}}
\def\QCoh{\mathit{QCoh}}
\def\Ker{\mathop{\mathrm{Ker}}}
\def\Im{\mathop{\mathrm{Im}}}
\def\Coker{\mathop{\mathrm{Coker}}}
\def\Coim{\mathop{\mathrm{Coim}}}

% Boxtimes
%
\DeclareMathSymbol{\boxtimes}{\mathbin}{AMSa}{"02}

%
% Macros for moduli stacks/spaces
%
\def\QCohstack{\mathcal{QC}\!\mathit{oh}}
\def\Cohstack{\mathcal{C}\!\mathit{oh}}
\def\Spacesstack{\mathcal{S}\!\mathit{paces}}
\def\Quotfunctor{\mathrm{Quot}}
\def\Hilbfunctor{\mathrm{Hilb}}
\def\Curvesstack{\mathcal{C}\!\mathit{urves}}
\def\Polarizedstack{\mathcal{P}\!\mathit{olarized}}
\def\Complexesstack{\mathcal{C}\!\mathit{omplexes}}
% \Pic is the operator that assigns to X its picard group, usage \Pic(X)
% \Picardstack_{X/B} denotes the Picard stack of X over B
% \Picardfunctor_{X/B} denotes the Picard functor of X over B
\def\Pic{\mathop{\mathrm{Pic}}\nolimits}
\def\Picardstack{\mathcal{P}\!\mathit{ic}}
\def\Picardfunctor{\mathrm{Pic}}
\def\Deformationcategory{\mathcal{D}\!\mathit{ef}}
\newenvironment{reference}{}{}
\newenvironment{slogan}{}{}
\newenvironment{history}{}{}
\usepackage{cleveref}
\crefname{originallemma}{original lemma}{original lemmas}
\Crefname{originallemma}{Original lemma}{Original lemmas}
\makeatletter
\newlabel{algebra-lemma-quotient-strict-henselization}{{algebra-lemma-quotient-strict-henselization}{source}{}{sourcepointer-1}{}}
\newlabel{algebra-lemma-quotient-strict-henselization@cref}{{[originallemma][0][]algebra-lemma-quotient-strict-henselization}{[1][0][]source}{}{}{}}
\makeatother
\crefname{originallemma}{original lemma}{original lemmas}
\Crefname{originallemma}{Original lemma}{Original lemmas}
\makeatletter
\newlabel{duality-lemma-gorenstein}{{duality-lemma-gorenstein}{source}{}{sourcepointer-2}{}}
\newlabel{duality-lemma-gorenstein@cref}{{[originallemma][0][]duality-lemma-gorenstein}{[1][0][]source}{}{}{}}
\makeatother
\crefname{originallemma}{original lemma}{original lemmas}
\Crefname{originallemma}{Original lemma}{Original lemmas}
\makeatletter
\newlabel{dualizing-lemma-completion-henselization-dualizing}{{dualizing-lemma-completion-henselization-dualizing}{source}{}{sourcepointer-3}{}}
\newlabel{dualizing-lemma-completion-henselization-dualizing@cref}{{[originallemma][0][]dualizing-lemma-completion-henselization-dualizing}{[1][0][]source}{}{}{}}
\makeatother
\crefname{originallemma}{original lemma}{original lemmas}
\Crefname{originallemma}{Original lemma}{Original lemmas}
\makeatletter
\newlabel{spaces-morphisms-lemma-flat-at-point-etale-local-rings}{{spaces-morphisms-lemma-flat-at-point-etale-local-rings}{source}{}{sourcepointer-4}{}}
\newlabel{spaces-morphisms-lemma-flat-at-point-etale-local-rings@cref}{{[originallemma][0][]spaces-morphisms-lemma-flat-at-point-etale-local-rings}{[1][0][]source}{}{}{}}
\makeatother
\crefname{originallemma}{original lemma}{original lemmas}
\Crefname{originallemma}{Original lemma}{Original lemmas}
\makeatletter
\newlabel{spaces-properties-lemma-describe-etale-local-ring}{{spaces-properties-lemma-describe-etale-local-ring}{source}{}{sourcepointer-5}{}}
\newlabel{spaces-properties-lemma-describe-etale-local-ring@cref}{{[originallemma][0][]spaces-properties-lemma-describe-etale-local-ring}{[1][0][]source}{}{}{}}
\makeatother
\title{Gorenstein morphisms of algebraic spaces tested on etale charts and at points via strict local rings}
\author{The Stacks Project Authors}
\date{stacks-project:spaces-more-morphisms.tex:local-snapshot-20261011; GNU FDL 1.2 or later}
\begin{document}
\maketitle
\noindent\textit{Editorial note.} Statement span is lines 6537-6560 of sources/stacks-spaces-more-morphisms/pinned-original-source.tex; proof span is lines 6562-6609, the complete author proof (46 body lines).\par\medskip
\noindent\textit{Editorial note.} Navigation stubs: 5 cited-but-not-included labels are supplied as non-load-bearing pointers with their pinned tags where the snapshot tag index has them.\par\medskip
\noindent\textit{Editorial note.} extra\_packages is empty on purpose: the retained preamble already loads xy with options and a second declaration clashes in the assembler wrapper.\par\medskip
\noindent\textit{Editorial note.} Text is retained verbatim, including the author's own typos and placeholders.\par\medskip
\noindent\textit{Allowed standard prerequisites (editorial).}\begin{itemize}
\item Explicit prerequisite boundary. Gorenstein morphisms of algebraic spaces, etale charts, strict local rings, and the schemes criterion for the Gorenstein property; the descent is retained.
\end{itemize}\medskip
\ifdefined\hypertarget\hypertarget{sourcepointer-1}{}\fi
\noindent\textit{Source reference.} Original lemma algebra-lemma-quotient-strict-henselization: sources/stacks-spaces-more-morphisms/pinned-original-source.tex, source lines 6609--6609. Cited by the retained statement or proof but not part of the collected spans; supplied only as a navigation pointer. The argument does not depend on its content.\par\medskip
\ifdefined\hypertarget\hypertarget{sourcepointer-2}{}\fi
\noindent\textit{Source reference.} Original lemma duality-lemma-gorenstein: sources/stacks-spaces-more-morphisms/pinned-original-source.tex, source lines 6609--6609. Cited by the retained statement or proof but not part of the collected spans; supplied only as a navigation pointer. The argument does not depend on its content.\par\medskip
\ifdefined\hypertarget\hypertarget{sourcepointer-3}{}\fi
\noindent\textit{Source reference.} Original lemma dualizing-lemma-completion-henselization-dualizing: sources/stacks-spaces-more-morphisms/pinned-original-source.tex, source lines 6609--6609. Cited by the retained statement or proof but not part of the collected spans; supplied only as a navigation pointer. The argument does not depend on its content.\par\medskip
\ifdefined\hypertarget\hypertarget{sourcepointer-4}{}\fi
\noindent\textit{Source reference.} Original lemma spaces-morphisms-lemma-flat-at-point-etale-local-rings: sources/stacks-spaces-more-morphisms/pinned-original-source.tex, source lines 6609--6609. Cited by the retained statement or proof but not part of the collected spans; supplied only as a navigation pointer. The argument does not depend on its content.\par\medskip
\ifdefined\hypertarget\hypertarget{sourcepointer-5}{}\fi
\noindent\textit{Source reference.} Original lemma spaces-properties-lemma-describe-etale-local-ring: sources/stacks-spaces-more-morphisms/pinned-original-source.tex, source lines 6609--6609. Cited by the retained statement or proof but not part of the collected spans; supplied only as a navigation pointer. The argument does not depend on its content.\par\medskip
\begin{lemma}
\label{lemma-gorenstein}
Let $S$ be a scheme. Let $f : X \to Y$ be a morphism of algebraic spaces
over $S$. Assume the fibres of $f$ are locally Noetherian.
The following are equivalent
\begin{enumerate}
\item $f$ is Gorenstein,
\item $f$ is flat and for some surjective \'etale morphism $V \to Y$
where $V$ is a scheme, the fibres of $X_V \to V$
are Gorenstein algebraic spaces, and
\item $f$ is flat and for any \'etale morphism $V \to Y$
where $V$ is a scheme, the fibres of $X_V \to V$
are Gorenstein algebraic spaces.
\end{enumerate}
Given $x \in |X|$ with image $y \in |Y|$ the following are
equivalent
\begin{enumerate}
\item[(a)] $f$ is Gorenstein at $x$, and
\item[(b)] $\mathcal{O}_{Y, \overline{y}} \to \mathcal{O}_{X, \overline{x}}$
is flat and
$\mathcal{O}_{X, \overline{x}}/
\mathfrak m_{\overline{y}}\mathcal{O}_{X, \overline{x}}$ is Gorenstein.
\end{enumerate}
\end{lemma}

\begin{proof}
Given an \'etale morphism $V \to Y$ where $V$ is a scheme
choose a scheme $U$ and a surjective \'etale morphism $U \to X \times_Y V$.
Consider the commutative diagram
$$
\xymatrix{
U \ar[d] \ar[r] & V \ar[d] \\
X \ar[r] & Y
}
$$
Let $u \in U$ with images $x \in |X|$, $y \in |Y|$, and $v \in V$.
Then $f$ is Gorenstein at $x$ if and only if $U \to V$ is
Gorenstein at $u$ (by definition). Moreover the morphism
$U_v \to X_v = (X_V)_v$ is surjective \'etale. Hence the scheme $U_v$ is
Gorenstein if and only if the algebraic space $X_v$ is Gorenstein.
Thus the equivalence of (1), (2), and (3) follows from the
corresponding equivalence for morphisms of
schemes, see Duality for Schemes, Lemma \ref{duality-lemma-gorenstein}
by a formal argument.

\medskip\noindent
Proof of equivalence of (a) and (b). The corresponding equivalence
for flatness is Morphisms of Spaces, Lemma
\ref{spaces-morphisms-lemma-flat-at-point-etale-local-rings}.
Thus we may assume $f$ is flat at $x$ when proving the equivalence.
Consider a diagram and $x, y, u, v$ as above. Then
$\mathcal{O}_{Y, \overline{y}} \to \mathcal{O}_{X, \overline{x}}$
is equal to the map
$\mathcal{O}_{V, v}^{sh} \to \mathcal{O}_{U, u}^{sh}$
on strict henselizations of local rings, see
Properties of Spaces, Lemma
\ref{spaces-properties-lemma-describe-etale-local-ring}.
Thus we have
$$
\mathcal{O}_{X, \overline{x}}/
\mathfrak m_{\overline{y}}\mathcal{O}_{X, \overline{x}} =
(\mathcal{O}_{U, u}/\mathfrak m_v \mathcal{O}_{U, u})^{sh}
$$
by Algebra, Lemma \ref{algebra-lemma-quotient-strict-henselization}.
Thus we have to show that the Noetherian local ring
$\mathcal{O}_{U, u}/\mathfrak m_v \mathcal{O}_{U, u}$
is Gorenstein if and only if its strict henselization is.
This follows immediately from
Dualizing Complexes, Lemma
\ref{dualizing-lemma-completion-henselization-dualizing}
and the definition of a Gorenstein local ring as a
Noetherian local ring which is a dualizing complex over itself.
\end{proof}

\end{document}
